% Part: turing-machines
% Chapter: undecidability
% Section: representing-tms

\documentclass[../../../include/open-logic-section]{subfiles}

\begin{document}

\olfileid{tur}{und}{rep}
\olsection{表示Turing机}

\begin{explain}
为了用一阶逻辑的语句来表示Turing机及其行为，我们必须定义一个合适的语言。
该语言由两部分组成：用于描述机器配置的谓词符号，以及用于对执行步骤（“时刻”）和纸带位置进行编号的表达式。

我们引入两类谓词符号，它们都是二元的：对于每个状态~\(q\)，引入一个谓词符号~\(\Obj Q_q\)；对于每个纸带符号~\(\sigma\)，引入一个谓词符号~\(\Obj S_\sigma\)。前者用于描述~\(M\) 的状态及其读写头的位置，后者用于描述纸带的内容。

为了表达读写头的位置和已执行的步数，我们需要一种表达数的方式。
这通过一个常项符号~\(\Obj 0\) 和一个一元函数~\(\prime\)（即后继函数）来实现。
按惯例，它写在自变元\emph{之后}（并且我们省略括号）。

对于每个数~\(n\)，都有一个标准项~\(\num{n}\)，即~\(n\) 的\emph{数词}，它在~\(\Lang L_M\) 中表示~\(n\)。
\(\num{0}\) 是 \(\Obj 0\)，\(\num{1}\) 是 \(\Obj 0'\)，\(\num{2}\) 是 \(\Obj 0''\)，依此类推。
更形式化地：
\begin{align*}
\num{0} & = \Obj 0 \\
\num{n+1} &= \num{n}'
\end{align*}
项 \(\num{0}\)（即 \(\Obj 0\)）命名纸带最左端的位置，以及第一个执行步骤之前的时刻（初始配置）。
项 \(\num{1}\)（即 \(\Obj 0'\)）命名最左端方格右侧的方格，以及第一个执行步骤之后的时刻，依此类推。

我们还引入一个谓词符号~\(<\)，用以表达纸带位置之间的顺序（此时意为“在……左边”）和执行步骤之间的顺序（此时意为“在……之前”）。

语言确定后，我们便可以列出 \(!T(M, w)\) 的“公理”，即那些共同描述~\(M\) 在输入~\(w\) 上运行行为的语句。
这些语句包括：为 \(\Obj 0\)、\(\prime\) 和 \(<\) 设定条件的语句，描述输入配置的语句，以及描述~\(M\) 在执行特定指令后的配置的语句。
\end{explain}

\begin{defn}
  \ollabel{defn:tm-descr}
给定一台Turing机 $M = \tuple{Q, \Sigma, q_0, \delta}$，其语言~\(\Lang L_M\) 由以下符号组成：
\begin{enumerate}
\item 对于每个状态~\(q \in Q\)，一个二元谓词符号 \(\Obj Q_q(x, y)\)。
  直观上，$\Obj Q_q(\num{m}, \num{n})$ 表示“在第 \(n\) 步之后，\(M\) 处于状态~\(q\) 并正在扫描第 \(m\) 个方格”；
\item 对于每个符号~\(\sigma\in \Sigma\)，一个二元谓词符号 \(\Obj S_\sigma(x, y)\)。
  直观上，$\Obj S_\sigma(\num{m},
  \num{n})$ 表示“在第 \(n\) 步之后，第 \(m\) 个方格包含符号~\(\sigma\)”；
\item 一个常项符号 \(\Obj 0\)；
\item 一个一元函数符号 \(\prime\)；
\item 一个二元谓词符号 \(<\)。
\end{enumerate}
\end{defn}

描述Turing机~\(M\) 在输入 $w = \sigma_{i_1}\dots\sigma_{i_k}$ 上运行的语句如下：
\begin{enumerate}
\item 描述数和~$<$ 的公理：
\begin{enumerate}
\item 表示每个数都小于其后继的语句：
\[
\lforall[x][x < x']
\]
\item 确保~$<$ 具有传递性的语句：
\[
\lforall[x][\lforall[y][\lforall[z][
      ((x < y \land y < z) \lif x < z)]]]
\]
\end{enumerate}
\item 描述输入配置的公理：
\begin{enumerate}
\item 在第~$0$ 步之后（即机器启动前），$M$ 处于初始状态~$q_0$，正在扫描方格~$1$：
\[
\Obj Q_{q_0}(\num{1}, \num{0})
\]
\item 前~$k+1$ 个方格包含符号 $\TMendtape$, $\sigma_{i_1}$, \dots, $\sigma_{i_k}$：
\[
\Obj S_\TMendtape(\num{0}, \num{0}) \land
\Obj S_{\sigma_{i_1}}(\num{1}, \num{0}) \land
\dots \land
\Obj S_{\sigma_{i_k}}(\num{k}, \num{0})
\]
\item 其余方格为空：
\[
\lforall[x][(\num{k} < x \lif \Obj S_\TMblank(x, \num{0}))]
\]
\end{enumerate}
\item 描述从一个配置转移到下一个配置的公理：

在下文中，令 $!A(x, y)$ 为所有形如
\[
\lforall[z][
  (((z < x \lor x < z) \land \Obj S_\sigma(z, y))
  \lif \Obj S_\sigma(z, y'))]
\]
（其中 $\sigma \in \Sigma$）的语句的合取。我们用 $!A(\num{m},\num{n})$ 来表示“除第~$m$ 个方格以外，第~$n+1$ 步之后的纸带与第~$n$ 步之后的纸带相同”。
\begin{enumerate}
\item \ollabel{rep-right} 对于每条指令 $\delta(q_i, \sigma) =
  \tuple{q_j, \sigma', \TMright}$，有语句：
\begin{align*}
& \lforall[x][\lforall[y][(
   (\Obj Q_{q_i}(x, y) \land \Obj S_{\sigma}(x, y)) \lif {}]] \\
&\qquad   (\Obj Q_{q_j}(x', y') \land \Obj S_{\sigma'}(x, y') \land
!A(x, y)))
\end{align*}
这条语句表示：如果在第~$y$ 步之后，机器处于状态~$q_i$ 并正在扫描包含符号~$\sigma$ 的方格~$x$，那么在第~$y+1$ 步之后，它将扫描方格~$x+1$，处于状态~$q_j$，方格~$x$ 现在包含~$\sigma'$，并且除~$x$ 以外的每个方格包含的符号与第~$y$ 步之后包含的相同。

\item \ollabel{rep-left} 对于每条指令 $\delta(q_i, \sigma) =
  \tuple{q_j, \sigma', \TMleft}$，有语句：
\begin{align*}
& \lforall[x][\lforall[y][
    ((\Obj Q_{q_i}(x', y) \land \Obj S_{\sigma}(x', y)) \lif {}]]\\
& \qquad   (\Obj Q_{q_j}(x, y') \land \Obj S_{\sigma'}(x', y') \land
!A(x, y))) \land {}\\
& \lforall[y][((\Obj Q_{q_i}(\num{0}, y) \land \Obj S_{\sigma}(\num{0},
    y)) \lif {}]\\
& \qquad (\Obj Q_{q_j}(\num{0}, y') \land \Obj S_{\sigma'}(\num{0},
  y') \land !A(\num{0}, y)))
\end{align*}
请思考一下这是如何运作的：现在我们不是以“如果正在扫描方格~$x$……”开头，而是：“如果正在扫描方格 $x+1$……”。向左移动意味着在下一步机器将扫描方格~$x$。但被写入的方格是~$x+1$。我们这样处理是因为我们没有减法或前驱函数。

注意，形如 $x+1$ 的数是 $1$, $2$, \dots，也就是说，这并未涵盖机器正在扫描方格~$0$ 且本应左移的情况（当然，它不能左移——它会停在原地）。这一特殊情况由第二个合取式涵盖：它表示，如果在第~$y$ 步之后，机器在状态~$q_i$ 下扫描方格~$0$ 且方格~$0$ 包含符号~$\sigma$，那么在第~$y+1$ 步之后，它仍然在扫描方格~$0$，现在处于状态~$q_j$，方格~$0$ 上的符号是 $\sigma'$，并且除方格~$0$ 以外的方格包含的符号与第~$y$ 步之后包含的相同。
\item \ollabel{rep-stay} 对于每条指令 $\delta(q_i, \sigma) =
  \tuple{q_j, \sigma', \TMstay}$，有语句：
\begin{align*}
& \lforall[x][\lforall[y][(
   (\Obj Q_{q_i}(x, y) \land \Obj S_{\sigma}(x, y)) \lif {}]] \\
&\qquad   (\Obj Q_{q_j}(x, y') \land \Obj S_{\sigma'}(x, y') \land
!A(x, y)))
\end{align*}
\end{enumerate}
\end{enumerate}

令 $!T(M, w)$ 为上述所有关于Turing机~$M$ 和输入~$w$ 的语句的合取。

为了表达~$M$ 最终会停机，我们必须找到一个表示“在某个步数之后，转移函数将无定义”的语句。令 $X$ 为所有使得~$\delta(q, \sigma)$ 无定义的序偶 $\tuple{q, \sigma}$ 的集合。令 $!E(M, w)$ 为如下语句：
\[
\lexists[x][\lexists[y][(\bigvee_{\tuple{q, \sigma} \in
      X}(\Obj Q_q(x, y) \land \Obj S_\sigma(x, y)))]]
\]

如果我们使用带有指定停机状态~$h$ 的Turing机，事情就更简单了：此时语句~$!E(M, w)$
\[
\lexists[x][\lexists[y][\Obj Q_h(x, y)]]
\]
就表示机器最终会停机。

\begin{prop}
\ollabel{prop:mlessk}
如果 $m < k$，则 $!T(M, w) \Entails \num{m} < \num{k}$
\end{prop}

\begin{proof}
练习。
\end{proof}

\begin{prob}
证明 \olref[tur][und][rep]{prop:mlessk}。
（提示：对 $k-m$ 使用归纳法）。
\end{prob}

\end{document}